%!TEX root = ../main.tex

\section{Соотношения баланса энергии}

{\nologo
\begin{frame}{Функционалы энергии}
\vspace{-5mm}
\begin{itemize}
    \item Плотность энергии электрического поля:
    \[
        w_{E} = \half \epsilon[\phi] \cdot \| \nabla \Phi \|_2^2.
    \]
    \item Плотность энергии, отнесенной к каналу пробоя:
    \[
        w_\text{br} = %
        \color{light_grey}\left[\color{black} \cfrac{\Gamma}{l^2} [1 - f(\phi)] \color{light_grey}\right]\color{black} + %
        \color{light_grey}\left[\color{black} \cfrac{\Gamma}{4} \| \nabla \phi \|_2^2 + \cfrac{\beta}{4} \Gamma l^2 \| \nabla \phi \|_2^4 \color{light_grey}\right]\color{black}.
    \]
\end{itemize}
\begin{columns}
\column{0.49\textwidth}
    \centering
    Свободная энергия Гельмгольца
    \[
        \psi = -w_E + w_\text{br}, \quad \Psi = \int\limits_\Omega \psi d \vx
    \]
\column{0.01\textwidth}
    \rule{0.4pt}{0.25\textheight}
\column{0.49\textwidth}
    \centering
    Внутренняя энергия
    \[
        w = w_E + w_\text{br}, \quad W = \int\limits_\Omega w d \vx
    \]
\end{columns}
\end{frame}
}


\begin{frame}{Балансовые соотношения}
\begin{block}{Утверждение}
    \begin{align*}
        \cfrac{d \Psi}{dt} &= -\cfrac{1}{m} \intOmega \left( \partt{\phi} \right)^2 d \vx + q^- \cfrac{d \Phi^+}{dt} - \intOmega \rho \partt{\Phi} d \vx; \\
        \cfrac{dW}{dt} &= -\cfrac{1}{m} \intOmega \left( \partt{\phi} \right)^2 d \vx - \Phi^+ \left( \cfrac{d q^-}{dt} + \int\limits_{y = 0} \sigma \cfrac{\partial \Phi}{\partial y} dS \right) - \intOmega Q_j d \vx.
    \end{align*}
    \vspace{-2mm}
\end{block}
\[
    Q_j = \sigma[\phi] \cdot (\nabla \Phi, \nabla \Phi).
\]
\end{frame}


{\nologo
\begin{frame}{Типы граничных условий}
\vspace{-6mm}
\begin{columns}
\column{0.3\textwidth}
    \vspace{0.85\textheight}
\column{0.01\textwidth}
    \rule{0.4pt}{0.85\textheight}
\column{0.3\textwidth}
    \vspace{0.85\textheight}
\column{0.01\textwidth}
    \rule{0.4pt}{0.85\textheight}
\column{0.3\textwidth}
    \vspace{0.85\textheight}
\end{columns}
\vspace{-0.83\textheight}
\begin{columns}
\column{0.33\textwidth}
    \centering
    Постоянный потенциал \\ на обкладках
\column{0.33\textwidth}
    \centering
    Постоянный заряд \\ на обкладках
\column{0.33\textwidth}
    \centering
    Сохранение заряда \\ в системе
\end{columns}
\vspace{2mm}
\noindent\rule[0.5ex]{\linewidth}{0.4pt}
\vspace*{-7mm}
\begin{columns}
\column{0.33\textwidth}
    \begin{gather*}
        \Phi^+ = \Const \\
        \sigma / \epsilon = \Const, \quad \rho_0 = 0 \\
        \rho \equiv 0
    \end{gather*}
\column{0.33\textwidth}
    \begin{gather*}
        q^- = -q^+ = \Const \\
        \sigma / \epsilon = \Const, \quad \rho_0 = 0 \\
        \rho \equiv 0
    \end{gather*}
\column{0.33\textwidth}
    \begin{gather*}
        \cfrac{d q^-}{dt} = -\int\limits_{y = 0} \sigma \cfrac{\partial \Phi}{\partial y} dS \\
        \cfrac{d q^+}{dt} = \int\limits_{y = H} \sigma \cfrac{\partial \Phi}{\partial y} dS
    \end{gather*}
\end{columns}
{\small
\begin{columns}
\column{0.33\textwidth}
    \[
        \cfrac{d \Psi}{dt} = -\cfrac{1}{m} \intOmega \left( \partt{\phi} \right)^2 d \vx
    \]
\column{0.33\textwidth}
    \[
        \cfrac{dW}{dt} = -\cfrac{1}{m} \intOmega \left( \partt{\phi} \right)^2 d \vx
    \]
\column{0.33\textwidth}
    \[
        \cfrac{dW}{dt} = -\cfrac{1}{m} \intOmega \left( \partt{\phi} \right)^2 d \vx - \intOmega Q_j d \vx.
    \]
\end{columns}
}
\vspace{3mm}
\begin{columns}
\column{0.33\textwidth}
    \centering
    $\Psi$ не возрастает
\column{0.33\textwidth}
    \centering
    $W$ не возрастает
\column{0.33\textwidth}
    \centering
    $W$ не возрастает
\end{columns}
\end{frame}
}